\section*{Capítulo \ref{ch:g1:animals-around-us} --- Los animales que nos rodean}

\begin{solution}{exo:g1:animals-around-us:1}
Cabeza, cuerpo, patas (\omterm{def:g1:animals-around-us:parts}{extremidades}) y cola. La cabeza lleva los ojos
y la boca.
\end{solution}

\begin{solution}{exo:g1:animals-around-us:2}
(a) $2$; (b) $4$; (c) $6$; (d) $8$; (e) ninguna --- la lombriz no
tiene patas.
\end{solution}

\begin{solution}{exo:g1:animals-around-us:3}
La oveja: lana, que es \omterm{def:g1:animals-around-us:covering}{pelo} espeso. El pato: \omterm{def:g1:animals-around-us:covering}{plumas}. La trucha:
escamas. El caracol: una \omterm{def:g1:animals-around-us:covering}{concha} (sobre un cuerpo blando y húmedo).
\end{solution}

\begin{solution}{exo:g1:animals-around-us:4}
El mirlo, la abeja y la mariposa. Todos tienen alas.
\end{solution}

\begin{solution}{exo:g1:animals-around-us:5}
El caracol se desliza despacio sobre su único pie blando y deja un
rastro de baba. La lombriz se retuerce: se estira larga y fina, luego
se encoge de nuevo, y avanza poco a poco.
\end{solution}

\begin{solution}{exo:g1:animals-around-us:6}
Anda con sus patas, nada con sus pies palmeados (otra vez las patas,
con los dedos unidos por \omterm{def:g1:animals-around-us:covering}{piel}) y vuela con sus alas.
\end{solution}

\begin{solution}{exo:g1:animals-around-us:7}
La hormiga se alimenta comiendo otros \omterm{def:g1:living-or-not:living}{seres vivos} (migas de comida,
semillas, animalillos) y se desplaza por sí misma con sus seis patas.
Alimentarse y moverse: encajan las dos partes de la definición.
\end{solution}

\begin{solution}{exo:g1:animals-around-us:8}
Respuesta abierta. Una buena descripción nombra las partes vistas, da
un número real de patas, la cubierta, la manera de moverse y el lugar
--- por ejemplo: ``un caracol en la pared: cabeza con dos tentáculos,
cuerpo blando, sin patas, \omterm{def:g1:animals-around-us:covering}{concha} marrón, se desliza; visto después de
la lluvia''.
\end{solution}

\begin{solution}{exo:g1:animals-around-us:9}
Andar es solo una manera de moverse. Los peces tienen aletas y nadan
--- esa es la suya. Un \omterm{def:g1:animals-around-us:animal}{animal} tiene que alimentarse y moverse, no
andar: el pez hace las dos cosas, así que es un \omterm{def:g1:animals-around-us:animal}{animal}.
\end{solution}

\begin{solution}{exo:g1:animals-around-us:10}
Sí, es un pájaro: las \omterm{def:g1:animals-around-us:covering}{plumas} y las alas lo delatan, aunque no vuele.
El avestruz no puede volar, pero corre más que cualquier persona ---
ser pájaro es cosa del cuerpo, no de volar.
\end{solution}
